\section{Related work}
\label{sec:related}

\paragraph{Evaluating LM personas.}
There has been growing interest in probing the ability of LMs to mimic human behaviors.
One line of work asks whether LMs can replicate results from well-known human experiments, \eg, in cognitive science, social science, and economics~\citep{uchendu2021turingbench,karra2022ai,aher2022using,binz2022using,srivastava2022beyond}.
Another set of studies have examined whether LMs can be used to simulate personas~\citep{park2022social,argyle2022out,jiang2022communitylm,simmons2022moral}, akin to our notion of steerability.
Through case studies in specific settings, these works gauge whether prompting LMs with demographic information (\eg, political identity) leads to human-like responses: \citet{argyle2022out} look at voting patterns and word associations, and \citet{simmons2022moral} consider moral biases.
By leveraging public opinion surveys, we are able to improve our understanding of LM steerability in three ways: (i)  \emph{breadth}: both in the range of different topics and steering groups, (ii) \emph{distributional view}: gauging whether LMs can match the spectrum of opinions of a group rather than its modal opinion, and (iii) \emph{measurability}: using metrics grounded in human response distributions.

Finally, recent works have examined the slants in the opinions of LMs---by prompting them with contentious propositions/questions generated by LMs~\citet{perez2022discovering} or from political tests~\citet{hartmann2023political}.
Similar to our work, they find that human-feedback trained models often exhibit a left-leaning, pro-environmental stance.
However, since our approach is based on public opinion surveys, we can go beyond the modal perspective taken by these works (comparing models to dominant viewpoints of specific groups, \eg, pro-immigration for liberals).
Instead, we can ask whether models reflect the varied perspectives held by real human subpopulations, and if they do so (consistently) across a range of topics.
We find that these two perspectives can often lead to different conclusions---\eg, \model{text-davinci-003} while very pro-liberal based on the modal view, does not capture liberal viewpoints in a nuanced and consistent manner according to our study. 


\paragraph{Subjectivity in evaluations.} 
There has been a long-standing push within the NLP community to consider the subjective and affective dimensions of language in evaluating models~\citep{alm2011subjective}.
Prior works show that for many tasks---from toxicity detection~\citep{Gordon21,Gordon22,davani2022dealing,Sap22,goyal2022your}, ethics judgements~\citep{lourie2021scruples}, and inference~\citep{pavlick2019inherent}---there is inherent variability in what different humans consider the ``correct answer''.
These studies serve as a motivation for our work, where we approach the problem of evaluating opinions expressed by LMs through the use of surveys.

\paragraph{Human-LM alignment.}
There is a growing body of work seeking to make LMs more human-aligned~\citep{askell2021general,ouyang2022training,glaese2022improving,bai2022constitutional}.
While these works recognize the subjectivity of the alignment problem, they do not focus on it---seeking instead to identify values to encode in models and building techniques to do so.
Our work looks instead delves deeper into the issue of subjectivity, asking \emph{who} are the humans that we are/should be aligning the models to?




\paragraph{Bias, toxicity, and truthfulness.}
There is a long line of work studying the bias and fairness of NLP systems~\citep{nadeem2020stereoset,Dhamala21,DeArteaga19,Brown20,eval-harness,srivastava2022beyond,liang2022holistic,xu2021bot,perez2022red,ganguli2022red}. These works study properties of LMs such as bias, toxicity, and truthfulness, focusing on flagging undesirable outcomes when the gold standard behavior is somewhat well-defined (\eg, don't use slurs). Our work takes a complementary perspective: evaluating LMs on inherently subjective questions taken from Pew Research. This allows us to gain quantitative insights into the representativeness of opinions expressed by LMs on contentious but important topics such as religion or privacy.

